%%
%% ACM sigconf template — DSAIT4050 IR Group 10 Project Report
%%
\documentclass[sigconf]{acmart}

\AtBeginDocument{%
  \providecommand\BibTeX{{Bib\TeX}}}

%% Course submission — no ACM rights needed
\setcopyright{none}
\acmConference[DSAIT4050 '25]{Information Retrieval}{Q3 2025}{Delft, Netherlands}

\begin{document}

\title[Cross-Script Phonetic Name Retrieval]{Cross-Script Phonetic Name Retrieval via Byte-Level Contrastive Embeddings}

\renewcommand{\shortauthors}{Jumle et al.}

\author{Vedant Vivek Jumle}
\email{v.v.jumle@student.tudelft.nl}
\affiliation{%
  \institution{Delft University of Technology}
}

\author{TODO: Author 2}
\email{TODO@student.tudelft.nl}
\affiliation{%
  \institution{Delft University of Technology}
}

\author{TODO: Author 3}
\email{TODO@student.tudelft.nl}
\affiliation{%
  \institution{Delft University of Technology}
}

\author{TODO: Author 4}
\email{TODO@student.tudelft.nl}
\affiliation{%
  \institution{Delft University of Technology}
}

%%
%% CCS Concepts
%%
\begin{CCSXML}
<ccs2012>
 <concept>
  <concept_id>10010147.10010178</concept_id>
  <concept_desc>Information systems~Retrieval models and ranking</concept_desc>
  <concept_significance>500</concept_significance>
 </concept>
 <concept>
  <concept_id>10010147.10010257.10010293.10010294</concept_id>
  <concept_desc>Computing methodologies~Neural networks</concept_desc>
  <concept_significance>300</concept_significance>
 </concept>
 <concept>
  <concept_id>10010147.10010178.10010179.10010182</concept_id>
  <concept_desc>Information systems~Similarity measures</concept_desc>
  <concept_significance>300</concept_significance>
 </concept>
</ccs2012>
\end{CCSXML}

\ccsdesc[500]{Information systems~Retrieval models and ranking}
\ccsdesc[300]{Computing methodologies~Neural networks}
\ccsdesc[300]{Information systems~Similarity measures}

\keywords{cross-script retrieval, phonetic embeddings, byte-level encoding, contrastive learning, name matching}

\begin{abstract}
Retrieving person names across different writing systems is a fundamental challenge in multilingual
information retrieval: a query like ``Schwarzenegger'' should match its Arabic, Russian, or Chinese
phonetic equivalents, despite sharing no characters. Existing approaches based on edit distance or
phonetic hashing fail entirely once the script boundary is crossed.

We present a byte-level transformer encoder trained with InfoNCE contrastive loss and
Approximate Nearest Neighbour hard-negative mining. The model operates directly on raw UTF-8
bytes---requiring no tokeniser, no language-specific preprocessing, and no script detection.
We construct a dataset of 119,040 named entities with 4.67 million positive pairs spanning
8 non-Latin scripts (Arabic, Russian, Chinese, Japanese, Hebrew, Hindi, Greek, Korean), generated
via LLM-based phonetic transliteration.

Our model achieves 0.775 MRR and 0.897 R@10 overall, reducing the script gap in R@10 from
0.93 (edit-distance baseline) to 0.098---a 10$\times$ reduction---while outperforming the best
rule-based baseline (transliterate, 0.565 MRR) by 37\%.
\end{abstract}

\maketitle

\section{Introduction}

Retrieving a person by name is deceptively hard when the query and the indexed record
differ in script. A sanctions screening system must match \textit{Владимир Путин} against
\textit{Vladimir Putin}; a hospital patient-matching system must link \textit{محمد} to
\textit{Mohamed} or \textit{Muhammad}. In each case the same individual is referred to,
yet the surface forms share no characters. Classical approaches---Levenshtein edit
distance~\cite{levenshtein1966binary,navarro2001guided}, Soundex, and BM25~\cite{robertson2009bm25}---operate
on raw character sequences and fail completely once a script boundary is crossed.

Recent work has shown that even strong multilingual dense retrievers degrade significantly
when queries are transliterated rather than written in their native script~\cite{chari2025lost}.
This \textit{script gap} is most severe for short name entities, where there is no surrounding
semantic context to compensate for surface-level mismatch. Rule-based transliteration pipelines
(e.g., ICU transliterate) partially close the gap, but apply a fixed canonical mapping that cannot
capture the phonetic variation introduced by different romanization conventions or LLM-generated
phonetic perturbations.

We address the specific problem of \textbf{cross-script phonetic name retrieval}: given a query
name in any script or phonetic form, retrieve the correct named entity from a corpus indexed by
its English name. We propose a lightweight transformer encoder trained from scratch on raw UTF-8
bytes~\cite{xue2022byt5,vaswani2017attention}. Each character in any writing system decomposes
deterministically into 1--4 bytes from a fixed 256-token vocabulary, eliminating out-of-vocabulary
tokens entirely and requiring no script detection or language-specific preprocessing. The encoder
is trained with InfoNCE contrastive loss~\cite{oord2018cpc} and ANCE-style hard-negative
mining~\cite{robinson2021hardneg}, pushing phonetically equivalent names together and
phonetically distinct names apart in embedding space.

To train and evaluate the model, we construct a dataset of 119,040 named entities with 4.67
million positive name pairs covering 8 non-Latin scripts (Arabic, Russian, Chinese, Japanese,
Hebrew, Hindi, Greek, Korean), generated via LLM-based phonetic transliteration and augmented
with Wikidata ground truth. Evaluation distinguishes three query types: \textit{phonetic} (Latin
variant of the English name), \textit{script} (direct transliteration into a non-Latin script),
and \textit{combined} (phonetic Latin variant transliterated into a non-Latin script).

Our main contributions are: (1) a controlled benchmark for cross-script and phonetically perturbed
name retrieval; (2) a compact byte-level contrastive encoder that achieves 0.897~R@10 and 0.775~MRR
overall, reducing the script gap in R@10 from 0.93 to 0.098---a 10$\times$ reduction over
edit-distance baselines; and (3) a systematic comparison of classical string-similarity baselines,
rule-based transliteration, and dense retrieval, with a FAISS index ablation characterising the
recall--latency tradeoff~\cite{johnson2021faiss,malkov2020hnsw}.

\section{Related Work}

\noindent\textbf{Fuzzy name matching.}
Classical name matching relies on character-level edit distance~\cite{levenshtein1966binary,navarro2001guided}
and phonetic codes such as Soundex and Metaphone, which encode names by their approximate English
pronunciation. BM25~\cite{robertson2009bm25} with character $n$-grams provides a sparse retrieval
baseline. These methods are effective within a single script but fail categorically across script
boundaries: a Latin query and an Arabic document share no character-level overlap regardless of
phonetic similarity.

\noindent\textbf{Script gap in neural IR.}
Chari et al.~\cite{chari2025lost} show that strong multilingual dense retrievers such as
BGE-M3~\cite{chen2024bge} degrade severely when queries are transliterated rather than written in
their native script. Their mitigation strategy fine-tunes existing models on mixtures of native and
Latinised text, requiring large pretrained multilingual backbones. Liu et al.~\cite{liu2024translico}
(TransliCo) address the same barrier by contrasting sentence representations with their Latin
transliterations during pretraining, improving cross-script alignment in general NLP tasks. Both
works treat the script gap as a fine-tuning or pretraining problem for sentence-level encoders.
We instead ask whether it can be solved \emph{from scratch}, for the specific case of short name
entities where no semantic context is available to compensate for surface mismatch.

\noindent\textbf{Contrastive dense retrieval.}
Karpukhin et al.~\cite{karpukhin2020dpr} (DPR) establish the dual-encoder paradigm, training
encoders with in-batch negatives to maximise query--document similarity. Gao et al.~\cite{gao2021simcse}
(SimCSE) demonstrate that hard negatives substantially improve contrastive embedding quality, and
Robinson et al.~\cite{robinson2021hardneg} provide grounding for hard-negative selection. The InfoNCE
objective~\cite{oord2018cpc} underlies our training loss. We adopt the contrastive retrieval framework
but train entirely from scratch on raw UTF-8 bytes~\cite{xue2022byt5}, without inheriting pretrained
representations---motivated by the incompatibility of subword vocabularies with cross-script phonetic
alignment.


\section{Method}


\section{Dataset}


\section{Experiments}


\section{Concluding Remarks}


\begin{acks}
Course project for DSAIT4050 Information Retrieval, Q3 2025, Delft University of Technology.
\end{acks}

\bibliographystyle{ACM-Reference-Format}
\bibliography{sample-base}

%%
%% The following sections are excluded from the 4-page limit.
%%

\section*{Link to Repository}

The code, dataset pipeline scripts, trained model checkpoint, and evaluation scripts are available at:
\url{https://github.com/TODO/ir-group10}

\section*{Self-Assessment of Contributions}

% TODO: each group member fills in their own paragraph.

\noindent\textbf{Vedant Vivek Jumle (5296196):} TODO

\noindent\textbf{TODO Author 2:} TODO

\noindent\textbf{TODO Author 3:} TODO

\noindent\textbf{TODO Author 4:} TODO

\section*{Declaration of Generative AI Usage}

This project used generative AI tools as writing aids and for code assistance.
Specifically, \textbf{Claude Sonnet 4.6} (Anthropic) was used to assist with writing
refinement---improving clarity, grammar, and structure of sections drafted by the authors.
The original intellectual contributions, research design, experimental methodology,
and interpretation of results are entirely our own.
All AI-assisted text was reviewed, revised, and approved by the authors before inclusion.

In accordance with the CEUR-WS GenAI Usage Taxonomy, our usage falls under
\emph{Writing Aid}---the tool refined prose written by the authors, but did not generate
novel scientific content or conclusions.

\end{document}
